\paragraph{D4b (Refuted universal lifting of the Albanese fibration).}
\emph{Existing-source check.}
Grothendieck's earlier first-order Igusa obstruction to a compatible
abelian Picard subgroup is already recorded in
\cite[Remark 3.4, p.~238]{DGJ_FGAVI1962}.
The exact lifting assertions addressed here were read in the published
2013 text of \cite{DGJ_Partsch2013}, pp.~102--106.
No correction to these assertions was found in the retrieved
publisher and author-source search; this limited search is not a
claim that no correction exists. No priority claim is made for the
obstruction. Its fully specified deformation is D4a.

The explicit deformation $X/R$ of D4a, with
$R=\overline{\mathbb F}_2[\epsilon]/(\epsilon^2)$, does not
admit a lifting of the Albanese elliptic fibration
\[
 f_0:X_0\longrightarrow Y_0=F_0/\langle P\rangle.
\]
Here a lifting means a morphism $f:X\to C$ over $R$, where
$C/R$ is a smooth proper deformation of $Y_0$, and with
special fibre identified with $f_0$. The special fibration is
Jacobian: its section is induced by $\{0\}\times F_0$ in the
cover $E_0\times F_0$.

Suppose such $f$ exists. The relative Picard scheme
$B=\operatorname{Pic}^0_{C/R}$ is an elliptic abelian scheme,
and pullback is a homomorphism of $R$-group schemes
\[
 f^*:B\longrightarrow\operatorname{Pic}_{X/R}.
\]
Its image lies in $\operatorname{Pic}^{\tau}_{X/R}$, since
pullback preserves algebraic triviality on geometric fibres.
Its special fibre is precisely the canonical inclusion
$Y_0^\vee=\operatorname{Pic}^{\rm var}(X_0)
\hookrightarrow\operatorname{Pic}^{\tau}(X_0)$, because $f_0$
is the Albanese map. This is the homomorphism ruled out by
D4a, for every possible deformation $B$. Thus $f$ cannot
exist. No assertion about the scheme-theoretic image of
$f^*$, or its flatness, enters this contradiction.

This refutes the literal universal lifting assertions in
\cite[Proposition 6.6, pp.~102--103; Theorem 6.7, p.~104]{DGJ_Partsch2013}:
the opened published text states that every deformation of
a Jacobian bielliptic surface extends its elliptic fibration,
and then claims the same for the Albanese fibration of every
bielliptic surface. The special fibre above is exactly a
Jacobian ordinary Igusa surface, and D4a constructs an actual
smooth projective deformation. The stronger assertion that
both fibrations lift, in Proposition~6.9, p.~106, is refuted
by the same example. This does not make a claim about other
cases of that article's classification.

The last trace step of the proof on p.~103 illustrates a
specific danger. Killing a connected subgroup and inducing
zero on an etale quotient does not force the original map to
vanish: its remaining map may go from that etale quotient to
a connected infinitesimal target. Over this very $R$ there
is a nontrivial homomorphism
\[
 h:(\mathbb Z/2)_R\longrightarrow\mu_{2,R},\qquad
 h(0)=1,\quad h(1)=1+\epsilon,
\]
whose special fibre is zero. The relation
$(1+\epsilon)^2=1$ verifies the group law. Thus rigidity of
maps between etale group schemes cannot be applied when
only the source is etale. The complete geometric refutation
above uses D4a; this finite-group example identifies an
invalid inference without being substituted for that proof.
The Hopf-map and group-law identities are recorded by
\texttt{check\_partsch\_etale\_step.py} and its saved output.
